% ============ 第2章 §2.2 Poisson核与Cauchy核(原书 pp.126-130) ============
\section{Poisson 核与 Cauchy 核}\label{sec:2.2}
\index{Poisson 核!(单位圆板的)}
\index{Cauchy 核}
\index{Poisson 积分}
\index{Cauchy 积分}
\index{Cauchy--Stieltjes 积分}
\index{Poisson--Stieltjes 积分}
\index{Herglotz 积分}
\index{Jensen 公式}
\index{零点集}


2.1 节中出现的 $(1-z)^{-1}$ 与 $\operatorname{Re}\left\{(1+z)(1-z)^{-1}\right\}$ 是非常重要的函数.$(1-z)^{-1}$ 称为 Cauchy 核,$\operatorname{Re}\left\{(1+z)(1-z)^{-1}\right\}$ 称为 Poisson 核.这两个核以如下方式相联系.对 $z \in D$,设 $z^{*}$ 为关于 $\partial D$ 的镜像,则 $z^{*} = z/|z|^{2}$,故
\begin{equation*}
\frac{1}{1 - z} - \frac{1}{1 - z^{*}} = \frac{1 - |z|^{2}}{|1 - z|^{2}} = \operatorname{Re}\left(\frac{1 + z}{1 - z}\right)
\end{equation*}
成立.用极坐标表示,设 $z = re^{i\theta}$($r \neq 1$),记
\begin{equation*}
\frac{1}{1 - re^{i\theta}} = C(r, \theta)
\end{equation*}
且
\begin{equation*}
\operatorname{Re}\left(\frac{1 + re^{i\theta}}{1 - re^{i\theta}}\right) = \frac{1 - r^{2}}{1 - 2r\cos\theta + r^{2}} = P(r, \theta).
\end{equation*}
若 $f$ 在包含 $\overline{D}$ 的开圆盘上正则,则由 Cauchy 积分表示与 Cauchy 积分定理
\begin{align*}
f(z) & = \frac{1}{2\pi i}\int_{|\zeta| = 1}\frac{f(\zeta)}{\zeta - z}\dd\zeta \\
& = \frac{1}{2\pi i}\int_{|\zeta| = 1}f(\zeta)\left(\frac{1}{\zeta - z} - \frac{1}{\zeta - z^{*}}\right)\dd\zeta
\end{align*}
因此
\begin{align*}
f(re^{i\theta}) & = \frac{1}{2\pi}\int_{-\pi}^{\pi}f(e^{it})C(r, \theta - t)\dd t \\
& = \frac{1}{2\pi}\int_{-\pi}^{\pi}f(e^{it})P(r, \theta - t)\dd t
\end{align*}
成立.

$M = M[-\pi, \pi]$ 表示 $[-\pi, \pi]$ 上(有界)复 Borel 测度的全体,对 $\mu \in M$ 令 $\|\mu\| = |\mu|([-\pi, \pi])$.对 $0 < p \leq \infty$,用 $L^{p} = L^{p}[-\pi, \pi]$ 表示 $[-\pi, \pi]$ 上 $p$ 次 Lebesgue 可积函数的全体,则 $L^{1}$ 可以等距嵌入 $M$.事实上,对 $F \in L^{1}$,令 $\mu = F(t)\dd t$,则 $\mu \in M$ 且 $\|\mu\| = \int_{-\pi}^{\pi}|F(t)|\dd t = \|F\|_{1}$.对 $\mu \in M$,
\begin{equation*}
f(re^{i\theta}) = \frac{1}{2\pi}\int_{-\pi}^{\pi}C(r, \theta - t)\dd\mu(t)
\end{equation*}
称为 Cauchy--Stieltjes 积分,
\begin{equation*}
u(re^{i\theta}) = \frac{1}{2\pi}\int_{-\pi}^{\pi}P(r, \theta - t)\dd\mu(t)
\end{equation*}
称为 Poisson--Stieltjes 积分.当 $\mu = F(t)\dd t$ 时,分别称为 Cauchy 积分,Poisson 积分.当 $\mu$ 是 Dirac 测度 $\delta_{t=0}$ 时,其 Cauchy--Stieltjes 积分是 $C(r, \theta)$,Poisson--Stieltjes 积分是 $P(r, \theta)$.

由于 $C(r, \theta - t) = e^{it}/(e^{it} - re^{i\theta})$ 关于 $z = re^{i\theta}$ 正则,Cauchy--Stieltjes 积分属于 $H(D)$.

由于 $P(r, \theta - t) = \operatorname{Re}(e^{it} + re^{i\theta})/(e^{it} - re^{i\theta})$ 关于 $z = re^{i\theta}$ 调和,当 $\mu$ 是实测度时 Poisson--Stieltjes 积分属于 $h(D)$.事实上,Poisson--Stieltjes 积分属于 $h^{1}$ 且 $\mu$ 的 Cauchy--Stieltjes 积分属于 $\bigcap_{0 < p < 1}H^{p}$ 这一点将在 2.3 节证明.

正如 2.1 节所注意到的,对 $u \in h(D)$,存在 $f \in H(D)$ 使得 $\operatorname{Re}(f) = u$.这可以由以下事实证明:当 $u$ 在包含 $\overline{D}$ 的开集上调和时,令
\begin{equation*}
f(z) = \frac{1}{2\pi}\int_{-\pi}^{\pi}\frac{e^{it} + z}{e^{it} - z}u(e^{it})\dd t \quad (z \in D),
\end{equation*}
则 $f$ 在 $D$ 上正则,两边取实部,
\begin{equation*}
\operatorname{Re}\left\{f(z)\right\} = \frac{1}{2\pi}\int_{-\pi}^{\pi}P(r, \theta - t)u(e^{it})\dd t \quad (z \in D)
\end{equation*}
且 $\operatorname{Re}\left\{f(z)\right\} = u(z)$($z \in D$).因此,当 $f$ 在 $\overline{D}$ 上正则且无零点时,$\operatorname{Re}\left\{\log f(z)\right\} = \log|f(z)|$($z \in \overline{D}$),故
\begin{equation*}
\log|f(z)| = \frac{1}{2\pi}\int_{-\pi}^{\pi}P(r, \theta - t)\log|f(e^{it})|\dd t \quad (z \in D).
\end{equation*}
令 $z = 0$,
\begin{equation*}
\log|f(0)| = \frac{1}{2\pi}\int_{-\pi}^{\pi}\log|f(e^{it})|\dd t.
\end{equation*}
将这个结果推广到 $f$ 在 $D$ 内有零点的情况.这称为 Jensen 公式.

设 $\{a_{j}\}$ 为 $f$ 在 $D$ 中的零点集.若为零点阶数为 $\ell$,则连续写出 $\ell$ 次.当在 $\overline{D}$ 上正则的 $f(z)$ 在 $z = a$ 处为零点时,除去该零点时不用 $f(z)/(z - a)$,而令
\begin{equation*}
h(z) = f(z)\Big/\frac{z - a}{1 - \bar{a}z},
\end{equation*}
则 $|h(e^{i\theta})| = |f(e^{i\theta})|$.这个想法经常被使用.

\begin{proposition}\label{prop:2.2.1}
$f$ 在 $\overline{D}$ 上正则,$f(0) \neq 0$ 且在 $\partial D$ 上无零点.设 $\{a_{j}\}$ 为 $f$ 在 $D$ 内的零点集,则 $\{a_{j}\}$ 为有限 $n$ 个,且
\begin{equation*}
\log|f(0)| = \frac{1}{2\pi}\int_{-\pi}^{\pi}\log|f(e^{i\theta})|\dd\theta + \sum_{j = 1}^{n}\log|a_{j}|.
\end{equation*}
\end{proposition}

\begin{proof}
由于 $f$ 在 $\overline{D}$ 上正则,由一致定理,$D$ 中的零点有限.除去 $f$ 在 $D$ 中的零点,令
\begin{equation*}
h(z) = f(z)\Bigg/\prod_{j = 1}^{n}\frac{z - a_{j}}{1 - \bar{a}_{j}z},
\end{equation*}
则如上所注意的,
\begin{equation*}
\log|h(0)| = \frac{1}{2\pi}\int_{-\pi}^{\pi}\log|h(e^{i\theta})|\dd\theta,
\end{equation*}
由于 $h(0) = f(0)/\prod_{j=1}^{n}(-a_{j})$ 且 $|h(e^{i\theta})| = |f(e^{i\theta})|$,本命题得证.
\end{proof}

\begin{proposition}\label{prop:2.2.2}
Poisson 核与 Cauchy 核
\begin{equation*}
P(r, \theta) = \frac{1 - r^{2}}{1 - 2r\cos\theta + r^{2}}, \quad C(r, \theta) = \frac{1}{1 - re^{i\theta}}
\end{equation*}
具有如下性质:

(1) $P(r, \theta) > 0$($0 \leq r < 1$),

(2) $P(r, \theta + 2\pi) = P(r, \theta)$($0 \leq r < 1$),

(3) $\frac{1}{2\pi}\int_{-\pi}^{\pi}P(r, \theta)\dd\theta = 1$($0 \leq r < 1$),

(4) $\operatorname{Re}\left\{C(r, \theta)\right\} = \frac{1 - r\cos\theta}{1 - 2r\cos\theta + r^{2}} > 0$,

(5) $C(r, \theta) - C\left(\frac{1}{r}, \theta\right) = P(r, \theta)$.
\end{proposition}

\begin{proof}
除 (3) 外都显然.(3) 对 $\operatorname{Re}\left(\frac{1 + z}{1 - z}\right)$ 用 Gauss 平均值定理来证明.
\end{proof}

命题~\ref{prop:2.2.1} 若把 $\overline{D}$ 换成 $\{z \in \mathbb{C};\ |z| \leq r\}$,则稍微改变形式后仍然成立,后面将多次用到.命题~\ref{prop:2.2.2} 是 Poisson 核的重要性质,后面将用到.

对 $\mu \in M[-\pi, \pi]$,
\begin{equation*}
f(z) = \frac{1}{2\pi}\int_{-\pi}^{\pi}\frac{e^{it} + z}{e^{it} - z}\dd\mu(t) \quad (z \in D)
\end{equation*}
称为 Herglotz 积分.当 $\mu$ 是实测度时,其实部是 Poisson--Stieltjes 积分.利用这个关系,证明 Poisson--Stieltjes 积分的唯一性.

\begin{proposition}\label{prop:2.2.3}
$\mu \in M[-\pi, \pi]$ 是实测度且满足 $\mu(\{-\pi\}) = \mu(\{\pi\})$.若 $\mu$ 的 Poisson--Stieltjes 积分在 $D$ 上为零,则 $\mu \equiv 0$.
\end{proposition}

\begin{proof}
由上面注意到的 Herglotz 积分与 Poisson 积分的关系,若
\begin{equation*}
u(z) = \int_{-\pi}^{\pi}P(r, \theta - t)\dd\mu(t) = 0 \quad (z \in D),
\end{equation*}
则存在实常数 $\gamma$ 使得
\begin{equation*}
\int_{-\pi}^{\pi}\frac{e^{it} + z}{e^{it} - z}\dd\mu(t) = i\gamma \quad (z \in D).
\end{equation*}
由于
\begin{equation*}
\frac{e^{it} + z}{e^{it} - z} = 1 + 2\sum_{n = 1}^{\infty}e^{-int}z^{n},
\end{equation*}
故 $\int_{-\pi}^{\pi}e^{-int}\dd\mu(t) = 0$($n = 0, \pm 1, \pm 2, \ldots$),由 Fourier 系数的唯一性,$\mu \equiv 0$.
\end{proof}

以后,当涉及 Poisson 积分的唯一性时,设 $[-\pi, \pi]$ 上的测度满足 $\mu(\{-\pi\}) = \mu(\{\pi\})$.另外,同样地,$\mu$ 也可以看作 $\partial D$ 上的测度.
